% Blueprint of the wNEAR verification. Every Lean definition, specification and
% theorem below is extracted by LeanArchitect from the `@[blueprint]`
% annotations in ../../{Model,Spec,Proofs}; run `lake build :blueprint` first.
\input{../../.lake/build/blueprint/library/Model}
\input{../../.lake/build/blueprint/library/Spec}
\input{../../.lake/build/blueprint/library/Proofs}
\input{../../.lake/build/blueprint/library/Blueprint}

\chapter{Introduction}

\section{The contract}

\texttt{wrap.near} is the wrapped NEAR fungible token (NEP-141), binary
\texttt{contracts/w-near/contract.wasm} with code hash
\texttt{DL2f5xmZ44cQRP6vBMt94rvseH5d5KKnAGis2WkTAsh1}, reproduced byte for
byte from \texttt{near/core-contracts} at commit \texttt{54f8267}, built with
Rust 1.51.0, near-sdk 3.1.0 and near-contract-standards 3.1.0. The contract
has fifteen exported methods: initialisation (\texttt{new}), wrapping
(\texttt{near\_deposit}, \texttt{near\_withdraw}), the NEP-141 core
(\texttt{ft\_transfer}, \texttt{ft\_transfer\_call},
\texttt{ft\_resolve\_transfer}, \texttt{ft\_total\_supply},
\texttt{ft\_balance\_of}, \texttt{ft\_metadata}) and NEP-145 storage
management (\texttt{storage\_deposit}, \texttt{storage\_withdraw},
\texttt{storage\_unregister}, \texttt{storage\_balance\_bounds},
\texttt{storage\_balance\_of}, plus the legacy
\texttt{storage\_minimum\_balance}).

\section{Structure of the verification}

The verification is organised in three layers.

\begin{enumerate}
\item \textbf{Model} (Chapter~\ref{ch:model}). A mathematical description of
  the contract state and of every method, written by hand from the Rust
  source. Each method is a function from a state and a host context to either
  a panic or an outcome (new state, promises, return value, logs). The
  mathematical model must be audited against the source; deployed conclusions
  also require audited runtime and adapter semantics and a refinement proof.
\item \textbf{Specifications} (Chapter~\ref{ch:spec}). Properties of the
  model stated as propositions, and their proofs (Chapter~\ref{ch:proofs}).
  They never mention WebAssembly.
\item \textbf{Refinement} (Chapter~\ref{ch:refinement}). The statement that
  the deployed binary, executed by the NEAR runtime in Lean, behaves exactly
  as the model, for every input, every amount of attached gas and every
  resource state. This is what transfers every proven property to the real
  contract. It is the generic \texttt{NearInterface.Refines}, instantiated
  with the near-lean runtime and the wNEAR contract model; its proof for
  \texttt{contract.wasm} remains open.
\end{enumerate}

A check mark on a definition records that its Lean statement exists; it does
not prove a proposition stated by that definition. Completed property proofs
appear as separate theorems in Chapter~\ref{ch:proofs}; deployed refinement
remains open.

\chapter{Model}
\label{ch:model}

\section{Types and state}

Amounts are natural numbers; the \texttt{u128} / \texttt{u64} ranges of the
on-chain representation are side conditions of well-formedness rather than
wrapping arithmetic, so every overflow check of the contract is an explicit
branch of the model.

\inputleanmodule{Model.Types}
\inputleanmodule{Model.State}

\section{Host context}

\inputleanmodule{Model.Env}

\section{Calls, outcomes and panics}

\inputleanmodule{Model.Calls}

\section{Methods}

One definition per exported method, in source order of
\texttt{contracts/w-near/source/w-near} and of the near-contract-standards
\texttt{FungibleToken} it delegates to. \texttt{step} adds the checks the
\texttt{\#[near\_bindgen]} wrapper generates (deposit, private,
initialisation) and dispatches.

\inputleanmodule{Model.Methods}

\section{On-chain layout}

\inputleanmodule{Model.Storage}

\chapter{Specifications}
\label{ch:spec}

\section{Invariant}

\inputleanmodule{Spec.Invariants}

\section{Authorization}

\inputleanmodule{Spec.Authorization}

\section{Conservation}

\inputleanmodule{Spec.Conservation}

\section{Availability}

\inputleanmodule{Spec.Availability}

\section{Views and initialisation}

\inputleanmodule{Spec.Views}

\section{Storage}

\inputleanmodule{Spec.Storage}

\chapter{Proofs}
\label{ch:proofs}

\inputleanmodule{Proofs.Views}
\inputleanmodule{Proofs.Invariants}
\inputleanmodule{Proofs.Authorization}
\inputleanmodule{Proofs.Conservation}
\inputleanmodule{Proofs.Availability}
\inputleanmodule{Proofs.Storage}

\chapter{Refinement}
\label{ch:refinement}

The refinement statement is generic and lives in \texttt{near-lean/interface}
(\texttt{NearInterface.Refines}). Its only hypotheses are that the storage
represents a model state and that the context is one the runtime builds. It
says that the execution never reports an interpreter limitation; that with at
least the model's gas bound attached it agrees with the model exactly (return
value, storage, balance, storage usage, logs, receipts, failure) and uses at
most that much gas; that with less it agrees or fails for gas with no effects;
and that input that is no call fails with no effects.

\section{The contract model}

The model reads the context as near-sdk 3.1.0 does: account ids carry over
(nearcore's and the SDK's account-id rules are proved to agree), arguments
are decoded from JSON as serde\_json and the derived \texttt{Deserialize}
accept them, and the callback result is parsed as a JSON \texttt{U128}. A
successful \cref{def:step} is followed by the runtime's checks: promise
deposits against the liquid balance, and the storage stake of the new state.
Return values and receipts are compared as bytes, panics by their exact
messages. The gas bound pays the longest possible trie path for every storage
key a call touches; that it holds is part of the obligation, and the adapter
tests check it on adversarial tries.

\inputleanmodule{Spec.Contract}

\section{The obligation}

\texttt{NearRuntime.Boundary} instantiates the runtime: one function-call
action of near-lean's VM followed by the receipt's storage-stake check, with
every failure rolling back storage, balance, usage and receipts. The theorem
below, with the binary embedded, is the formal deployed refinement
obligation. It is not proved.

\inputleanmodule{Proofs.Refinement}

\section{What the statement buys}

\cref{thm:withdraw-available-binary} shows the use: under the refinement,
\cref{spec:withdraw-available-boundary} becomes a statement about the
execution of the binary itself.


\chapter{Current dependency frontier}
\label{ch:frontier}

The generated \href{frontier.html}{live dependency frontier} distinguishes
Proven (kernel checked under reviewed environmental premises), Conditional
(proved composition with unresolved verification obligations), Open (a precise
contract or admitted leaf without its proof), and Integrated (used by a checked
caller composition). Integration does not make a conditional caller complete.
Definitions state contracts; their formalization check marks do not discharge them.
Solid graph edges come from LeanArchitect's Lean dependency traversal. Planned
edges are separately labeled and are not evidence of integration. The frontier
records the source revision, source digest, last successful build, and failures.

\inputleanmodule{Blueprint}
